\documentclass[xcolor=pdftex,dvipsnames,table,notheorems]{beamer}
\usetheme{Madrid}
\usepackage{lmodern}
\setbeamercovered{dynamic}
\usepackage{wrapfig}
\setbeamertemplate{navigation symbols}{}
\usepackage[absolute,overlay]{textpos}
%\usepackage{hyperref}
\usepackage{graphics}
\usepackage[T1]{fontenc}
\usepackage{amsmath, amssymb,mathtools}
\usepackage[ansinew]{inputenc}
\usepackage{hyperref}
\usepackage{bbm} 
\usepackage{graphicx}
\usepackage{appendixnumberbeamer}
\usepackage{color}
%\usepackage{xcolor}
\usepackage{subcaption}
\usepackage{centernot}
\usepackage{booktabs}
\usepackage{threeparttable}
\usepackage{tikz}
\usetikzlibrary{arrows.meta}
\usepackage{pgfplots}
\pgfplotsset{compat=1.18}
\usepackage{comment}
\setbeamercovered{invisible}
\usepackage[dvipsnames]{xcolor}
\usepackage{soul}
\usepackage{xcolor}
\usepackage{comment}

\newcommand\independent{\protect\mathpalette{\protect\independenT}{\perp}}
\def\independenT#1#2{\mathrel{\rlap{$#1#2$}\mkern2mu{#1#2}}}

\usepackage{amsthm}
\setbeamertemplate{theorems}[numbered]
\newtheorem{theorem}{Theorem}
\newtheorem{corollary}{Corollary}[theorem]
\newtheorem{lemma}{Lemma}
\newtheorem{proposition}[theorem]{Proposition}
\newtheorem{assumption}{Assumption}
\newtheorem{definition}{Definition}


\newcommand*{\boxedcolor}{red}
\makeatletter
\renewcommand{\boxed}[1]{\textcolor{\boxedcolor}{%
  \fbox{\normalcolor\m@th$\displaystyle#1$}}}
  
  \newcommand{\sym}[1]{\rlap{#1}}

\let\estinput=\input% define a new input command so that we can still flatten the document

\newcommand{\estwide}[3]{
		\vspace{.85ex}{
			\begin{tabular*}
			{\textwidth}{@{\hskip\tabcolsep\extracolsep\fill}l*{#2}{#3}}
			\toprule
			\estinput{#1}
			\bottomrule
			\addlinespace[.85ex]
			\end{tabular*}
			}
		}

\newcommand\norm[1]{\left\lVert#1\right\rVert}


\captionsetup[figtext]{justification=justified,font=footnotesize}

\newcommand{\figtext}[1]{%
% 	\vspace{-1.9ex}
\captionsetup{options=figtext}
\caption*{\hspace{6pt}\hangindent=1.5em #1}
}
\newcommand{\Figtext}[1]{%
\begin{tablenotes}[para,flushleft]
\hspace{6pt}
\hangindent=1.85em
#1
\end{tablenotes}
}
\newcommand{\Fignote}[1]{\Figtext{\emph{Note:~}~#1}}
\newcommand{\Figsource}[1]{\Figtext{\emph{Source:~}~#1}}


\usecolortheme[named=Blue]{structure}

\setbeamertemplate{caption}[numbered]{}


\newcommand\bfootnote[1]{%
  \begingroup
  \renewcommand\thefootnote{}\footnote{#1}%
  \addtocounter{footnote}{-1}%
  \endgroup
}

\definecolor{col_before}{HTML}{489c74}
\definecolor{col_after}{HTML}{cc640c}

\newcommand{\roadmapdot}[2]{%
  \ifnum#1<#2
    \fill[col_before] (#1,0) circle (0.105);
  \else
    \ifnum#1=#2
      \fill[col_after] (#1,0) circle (0.13);
    \else
      \fill[black!18] (#1,0) circle (0.105);
    \fi
  \fi
}

\newcommand{\roadmapdots}[1]{%
  \begin{tikzpicture}[x=0.85cm,y=0.85cm,baseline]
    \foreach \i in {1,...,6} {\roadmapdot{\i}{#1}}
  \end{tikzpicture}
}

\newcommand{\roadmapitem}[3]{%
  \ifnum#2=#1
    \item \textbf{\textcolor{col_after}{#3}}
  \else
    \item #3
  \fi
}

\newcommand{\roadmapframe}[1]{%
  \begin{frame}{Roadmap}
    \centering
    \roadmapdots{#1}
    \vspace{0.8em}
    \begin{enumerate}
      \roadmapitem{#1}{1}{Institutional background and data}
      \roadmapitem{#1}{2}{Model}
      \roadmapitem{#1}{3}{Empirical strategy}
      \roadmapitem{#1}{4}{Welfare effect of the reform}
      \roadmapitem{#1}{5}{Disentangling level and slope responses}
      \roadmapitem{#1}{6}{Optimal reform}
    \end{enumerate}
  \end{frame}
}
